\section{An alias written where a lexer goes}
\label{sec:ch01-minted-alias}

This book declares \code{fixturesdl} with a \code{newminted} line in its
preamble, so it is an environment name. Written where the lexer goes it
reaches Pygments, which has never heard of it, and the build dies naming a
lexer problem the book does not have.

\begin{minted}{fixturesdl}
type Session { id: Int! }
\end{minted}

An option list does not hide the argument, so this is the second finding.

\begin{minted}[breaklines]{fixturesdl}
type Speaker { id: Int! }
\end{minted}

A block that opens and closes on one line fails the build the same way, so it
is the third: \begin{minted}{fixturesdl}x\end{minted}

The three below are not findings, and each is a different way of being right.

\begin{fixturesdl}
type Session { id: Int! }
\end{fixturesdl}

\begin{minted}{ruby}
puts 'the real lexer the alias borrows'
\end{minted}

\begin{minted}[fontsize=\footnotesize]{text}
an option list before a lexer that is not an alias
\end{minted}
